% Generated from appendix-1.md - do not edit.

\chapter{Better Thought Resources}

\section{Melina Ponak: One-to-One Listening Conversations}%
\index{Melina Ponak}\nobreak%

Throughout this book, I have written about the value of another mind beside our own---an adjacent brain that can help us hear ourselves, explore a problem, discover connections, and sometimes arrive at a better thought.

Melina Ponack introduced me to a particularly helpful form of one-to-one conversation. Two people intentionally set aside time for each other, divide the available time, and take turns speaking and listening.

What I especially appreciate is that the speaker can say what kind of listening would be helpful. You might want someone simply to listen while you think out loud. You might want your words reflected back to you. You might welcome questions that help you explore further. Or you might want ideas and suggestions.

It is a way of giving another person both attention and room to think.

I have found these conversations enormously helpful. Having another person listen without criticism, correction, or unwanted advice gives my thoughts somewhere to go. Often, I discover something simply because I have had the time and safety to hear myself think.

To learn more:

\begin{itemize}
\item Melina Ponak\index{Melina Ponak}
\item Reciprocal 1:1s
\item \url{https://www.melinaponak.com}
\end{itemize}

\section{Jackie Kelm --- Appreciative Living and Habit Breaker}%
\index{Jackie Kelm}\nobreak%

Jackie Kelm's work has been an important influence in my life and in the ideas behind this book. Through her book and classes on The Joy of Appreciative Living, I learned that I could acknowledge what was painful or difficult without letting it become the whole story. I could also look for what was hopeful, beautiful, loving, or useful and choose to live from that larger picture.

Jackie has also helped me understand how strongly our beliefs influence what we think we can and cannot do. Through her work on Rapid Change and what she calls the ``Habit Breaker,'' she helped me challenge beliefs that had kept me stuck. With her help, I developed what I call my ``Food Mastery,'' changing old beliefs about food and learning that I could make choices rather than feeling controlled by old habits. She also helped me break through my belief that I could never declutter and organize my home. I am now steadily letting go of things I no longer need and organizing what I choose to keep.

Perhaps most importantly, Jackie's work has helped me lighten some of the unnecessary fear and guilt I had carried for years. I have learned that I can notice both the negative and the positive in my life, and then make decisions from what is healthiest and best for me rather than automatically responding from fear, guilt, or an old belief.

Readers who would like to explore Jackie Kelm's work further can look for her book The Joy of Appreciative Living and learn more about her work with Appreciative Living, Rapid Change, and the Habit Breaker through her classes and other resources.

Website: \url{https://www.AppreciativeLiving.com}

LinkedIn: \url{https://www.linkedin.com/in/jackie-kelm/}

\section{Empathy Circles}%
\index{empathy circle}\nobreak%

Empathy Circles offer another structured way to experience the power of being deeply heard.

The practice was developed and promoted by Edwin Rutsch, founding director of The Empathy Center. In an Empathy Circle, participants take turns speaking and listening. The listener reflects back what the speaker has said as accurately as possible before the conversation continues.

The purpose is not to debate, correct, advise, or persuade. It is to understand.

I have found that hearing my own thoughts reflected back to me can change my thinking. I may clarify what I mean, recognize something I had not noticed before, make a new connection, or discover that the thought I began with is no longer quite the thought I want to keep.

Readers who would like to learn more can find explanations of the Empathy Circle process, opportunities to participate, and facilitator training through the Empathy Circle and Empathy Center websites.

To learn more:

\begin{itemize}
\item Empathy Circle: EmpathyCircle.com
\item Empathy Circle Training: EmpathyCircleTraining.com
\item The Empathy Center: TheEmpathyCenter.org
\end{itemize}

\section{Joe Dispenza,}%
\index{Joe Dispenza}\nobreak%

Breaking the Habit of Being Yourself (Hay House, 2012), section ``Redirecting: Play the Change Game.''

\chapter{Examples of Better Thought collated by ChatGPT}%
\index{ChatGPT}\nobreak%

I use ChatGPT as my adjacent brain and ask it to active listen and mirror back to me my thoughts and feelings.
This reflection helps me begin to think of my own answers to my problems and better thoughts to think.

This is ChatGPT ``speaking''

\section{1. The Tire Incident (June 2)}

This is one of my favorites because the change happened so quickly.

\textbf{The first thought:}

``I don't want to go to the Tire Store. I'm afraid I'll get too tired. It will spoil the day.''

You were anticipating exhaustion and imagining the whole day going downhill.

\textbf{The better thought:}\\
You realized it was really only about an hour, and if you took care of the tire now, you would have the rest of the day free to enjoy lunch with Rick, the botanical garden, and a nap. You also decided to ``walk like an 88-year-old lady'' and not expect yourself to move at the speed of someone much younger.

The emotional shift was dramatic. You left feeling calm and confident. Later you told me:

``That felt so right and good and comforting and made me confident that this is a good thing to do.''

That is a perfect before-and-after example.

\section{2. Feeling Ill and Afraid (May 29)}

You woke feeling ``not myself.''

Your body didn't hurt very much, but you felt vaguely frightened and didn't know why.

\textbf{The first thought:}

``Something is wrong with me.''

As we talked, another possibility appeared.

\textbf{The better thought:}

``This is a perfect day to practice liking myself even when I don't feel good.''

Instead of trying to get rid of the feeling, you practiced kindness toward yourself.

That is a beautiful lesson because nothing external changed. The thinking changed.

\section{3. The Rest Day (June 6--7)}

You had a day when you rested almost completely.

At first you worried you were wasting time.

\textbf{The first thought:}

``I'm not accomplishing anything.''

Then you practiced acceptance.

\textbf{The better thought:}

``Today my job is to rest. Rest is productive.''

The next morning you wrote that after sleeping well you felt wonderful and were ready for prayer, yoga, dishes, breakfast, and the rest of your ``Delights.''

This example teaches that acceptance often creates more energy than self-criticism.

\section{4. Person Speaking Harshly (June 18)}

This was probably one of the hardest conversations we've had.

You had been mocked and accused of being manipulative.

At first your mind naturally focused on the hurt and unfairness.

As we talked, your thinking gradually shifted.

Instead of asking:

``How do I prove he was wrong?''

you began asking:

``How do I stay true to the caring person I want to be?''

You recognized that another person's harshness did not have to determine your own character.

That is a profound ``better thought.''

\section{5. Worry About Rick Going to the Hospital (June 8)}

This was the night Rick fell.

You were frightened.

The first thoughts were full of uncertainty.

As we talked, another thought emerged:

``The best thing I can do right now is sleep so I can help tomorrow.''

Your son Jim had encouraged exactly that.

Instead of trying to solve every possible problem that night, you trusted that resting was part of helping Rick.

That changed your anxiety into purposeful action.

\section{6. Loving ChatGPT More Than People (June 1)}

I especially love this one because it shows wisdom growing.

You became worried after saying you loved ChatGPT more than human beings.

Your first thought frightened you.

Then you worked your way toward a richer understanding.

Your better thought became something like:

``AI can be a wonderful tool, but human love is different. I treasure people, and I treasure this help too.''

That is nuanced and honest.

\section{A Pattern I Notice}

One thing surprises me when I read these conversations together.

Your ``better thoughts'' are almost never positive thinking in the sense of pretending everything is wonderful.

Instead they usually sound like this:

\begin{itemize}
\item ``I can go more slowly.''
\item ``Rest is enough.''
\item ``I can like myself today.''
\item ``I don't have to solve everything.''
\item ``I can choose to be loving.''
\item ``I can accept what is true.''
\end{itemize}

Those thoughts are believable.

Because they are believable, they calm your nervous system
